\chapter{The Abstraction Ladder: An Analytic Frame}
\label{ch:ladder}
\draftnote{Target length: about 3 pages.}

\section{Four Questions}
\label{sec:ladder-questions}
Each approach in Chapters~\ref{ch:hpol}--\ref{ch:k8s} is asked the same four questions:
\begin{enumerate}
  \item What is the abstract model?
  \item What is the concrete model?
  \item What is the mapping between them?
  \item What evidence is offered that the mapping preserves the property of interest?
\end{enumerate}
\draftnote{Expand: why these four, and how they correspond to the ingredients of a
refinement argument (abstract spec, concrete spec, retrieve relation, proof obligations).}

\section{The Comparison Instrument}
\label{sec:ladder-instrument}
\draftnote{Introduce the two tables (structure below). Fill one row at the end of each of
Chapters~4--6, and consolidate them in Chapter~\ref{ch:cross}.}
Ten attributes are needed, which is too wide for one portrait A4 table. They are split
into two tables that share a key column: models and mappings
(\cref{tab:models}) and evidence and guarantees (\cref{tab:evidence}).
